% ======================================================================
%  PAPER II of the two-paper split
%  Brachistochrones in conformal Kerr spacetimes:
%  ergosphere trichotomy, separatrices, and adiabatic response
%  IOP iopart class -- Classical and Quantum Gravity
%  (axisymmetric sector: Thakurta-Kerr; control-theory foundations are
%   cited from Paper I)
% ======================================================================
\documentclass[12pt]{iopart}

% iopart pre-defines \equation*, which clashes with amsmath; free it first
\makeatletter
\expandafter\let\csname equation*\endcsname\relax
\expandafter\let\csname endequation*\endcsname\relax
\makeatother
\usepackage{amsmath}
\usepackage{iopams}
\usepackage{graphicx}
\usepackage{bm}
\usepackage{xcolor}
\usepackage{hyperref}
\hypersetup{colorlinks=true, linkcolor=blue, citecolor=blue, urlcolor=blue}

% figures shared with the master/Paper I
\graphicspath{{../paper/}}

% prose placeholder helper
\newcommand{\todo}[1]{{\color{red}\small[[#1]]}}

% shortcuts
\newcommand{\Ehat}{\hat E}
\newcommand{\rp}{r_{+}}
\newcommand{\re}{r_{e}}
\newcommand{\Jc}{J_{c}}
\newcommand{\Af}{A_{\mathrm{freeze}}}
\newcommand{\dd}{\mathrm{d}}

\begin{document}

\title{Brachistochrones in conformal Kerr spacetimes:
ergosphere trichotomy, separatrices, and adiabatic response}

\author{Iman Rosignoli}
\address{University of Pavia, Pavia, Italy}
\ead{iman.rosignoli01@universitadipavia.it}

\begin{abstract}
\todo{Paper II abstract (<=300 words, CQG): Thakurta-Kerr = conformal Kerr;
the breathing indicatrix and the arrival branch t==eta vs proper time tau;
frame-dragging; closed-form Weierstrass/Kleinian separatrices continued through
the ergosphere with the cusp theorem and trichotomy; rotational and conformal
plunge inversion with same-launch vs fixed-endpoint (tiered R\_*(E,a));
first-order on-shell/off-shell correction canonically re-derived and validated
against the true flow; length-two iterated Abelian integrals (weight-two,
irreducibility conjectural). Control-theory foundations are established in
Paper I.}
\end{abstract}

\maketitle
{\hypersetup{linkcolor=black}\tableofcontents}

% ======================================================================
\section{Introduction}
\todo{Adapt master intro; scope to the axisymmetric (Thakurta-Kerr) case;
cite Paper I for the controlled-rail formulation, existence/normality/HJB, and
the W-hierarchy.}

\section{Thakurta--Kerr: conformal rotating compact object}
\todo{Migrate master TK section: indicatrix and Hamiltonians (H\_eta, H\_t,
H\_tau); the conformal-time (eta) brachistochrone and t==eta arrival family;
Randers optical reduction; quasi-constants drift.
DO NOT migrate Sec. 5.1 (off-equatorial quasi-constants) -- hold for a later
3D paper.
Then: semi-analytic first-order adiabatic orbit shape; the branch costate vs
mechanical angular momentum dictionary.}

\section{Equatorial closed forms and the ergosphere}
\todo{Migrate master Sec.: separatrix in Weierstrass functions; trichotomy and
the ergosphere cusp; plunge inversion (rotational and conformal); fixed-endpoint
comparison and protocol dependence (tiered R\_*(E,a)); the breathing indicatrix
unified adiabatic theorem.}

\section{Conclusions}
\todo{Three blocks: Proved / Conditional and numerical / Outlook.}

\appendix
\section{Validation and consistency checks (Thakurta--Kerr)}
\todo{TK parts of the master validation appendix.}

\section{Explicit adiabatic forms on the genus-degeneration loci (Thakurta--Kerr)}
\todo{TK tau/t separatrices: physical external separatrix (r\_d>2M) vs algebraic
J\_deg roots.}

\section{Complete first-order extended-Hamiltonian correction (Thakurta--Kerr)}
\todo{Canonical derivation in (r,p\_r); conformal source S\_D=eta+int E H\_E d.eta;
length-two iterated Abelian integrals with weight-one/weight-two hierarchy.}

% ---- back matter -----------------------------------------------------
\section*{Data availability}
\todo{Same Zenodo DOI 10.5281/zenodo.21707378 statement as the master.}

\ack
\todo{Same acknowledgements as master.}

\nocite{*}
\bibliographystyle{iopart-num}
\bibliography{../paper/refs}

\end{document}
